\documentclass[12pt]{article} 

\usepackage{amsmath} 
\usepackage{amssymb}
\usepackage{tikz}
\usepackage{tikz-3dplot}
\usepackage[margin=0.6in]{geometry} 

\setlength{\abovedisplayskip}{4pt} 
\setlength{\belowdisplayskip}{4pt} 
\setlength{\abovedisplayshortskip}{2pt} 
\setlength{\belowdisplayshortskip}{2pt} 
\setlength{\parskip}{2pt} 
\setlength{\parsep}{0pt} 

\begin{document} 
	
	\begin{center} 
		{\large\textbf{APM1220}} \quad 
		\textbf{FA 2 -- Foundations: Sample Geometry} \quad  
		\\ Joseph Raphael Espiritu 
	\end{center} 
	
	\vspace{-0.3em} 
	
	\section*{Problem 3.1}
	
	Given the data matrix
	\[
	X=
	\begin{bmatrix}
		9&1\\
		5&3\\
		1&2
	\end{bmatrix}.
	\]
	
	There are \(n=3\) observations and \(p=2\) variables.
	
	\subsection*{(a)}
	
	The rows of \(X\) represent the three observations:
	\[
	\mathbf{x}_1'=(9,1),\qquad
	\mathbf{x}_2'=(5,3),\qquad
	\mathbf{x}_3'=(1,2).
	\]
	
	Therefore, the three points to be plotted in the
	\(p=2\) dimensional scatter plot are
	\[
	(9,1),\qquad (5,3),\qquad (1,2).
	\]
	
	The sample mean vector is
	\[
	\bar{\mathbf{x}}
	=
	\begin{bmatrix}
		\bar{x}_1\\
		\bar{x}_2
	\end{bmatrix}
	=
	\begin{bmatrix}
		\dfrac{9+5+1}{3}\\[6pt]
		\dfrac{1+3+2}{3}
	\end{bmatrix}
	=
	\begin{bmatrix}
		\dfrac{15}{3}\\[6pt]
		\dfrac{6}{3}
	\end{bmatrix}
	=
	\boxed{
		\begin{bmatrix}
			5\\
			2
	\end{bmatrix}}.
	\]
	
	Hence, the sample mean is located at
	\[
	\boxed{\bar{\mathbf{x}}=(5,2)}.
	\]
	
	The scatter plot with the sample mean located on the diagram
	is shown below.
	
	\begin{center}
		\begin{tikzpicture}[scale=0.65]
			
			% Axes
			\draw[->] (0,0) -- (10.5,0) node[right] {$x_1$};
			\draw[->] (0,0) -- (0,4.5) node[above] {$x_2$};
			
			% x-axis ticks
			\foreach \x in {1,...,10}
			{
				\draw (\x,0.08) -- (\x,-0.08);
				\node[below] at (\x,-0.08) {\small \x};
			}
			
			% y-axis ticks
			\foreach \y in {1,...,4}
			{
				\draw (0.08,\y) -- (-0.08,\y);
				\node[left] at (-0.08,\y) {\small \y};
			}
			
			% Data points
			\fill (9,1) circle (2.5pt);
			\node[above right] at (9,1) {$(9,1)$};
			
			\fill (5,3) circle (2.5pt);
			\node[above right] at (5,3) {$(5,3)$};
			
			\fill (1,2) circle (2.5pt);
			\node[above right] at (1,2) {$(1,2)$};
			
			% Sample mean
			\draw[thick] (4.85,1.85) -- (5.15,2.15);
			\draw[thick] (4.85,2.15) -- (5.15,1.85);
			\node[below right] at (5,2) {$\bar{\mathbf{x}}=(5,2)$};
			
		\end{tikzpicture}
	\end{center}
	
	\subsection*{(b)}
	
	For the \(n=3\)-dimensional representation, the columns of
	\(X\) are treated as vectors in \(\mathbb{R}^3\).
	
	The first column is
	\[
	\mathbf{y}_1
	=
	\begin{bmatrix}
		9\\
		5\\
		1
	\end{bmatrix},
	\]
	and the second column is
	\[
	\mathbf{y}_2
	=
	\begin{bmatrix}
		1\\
		3\\
		2
	\end{bmatrix}.
	\]
	
	Let
	\[
	\mathbf{1}
	=
	\begin{bmatrix}
		1\\
		1\\
		1
	\end{bmatrix}.
	\]
	
	From part (a),
	\[
	\bar{x}_1=5,
	\qquad
	\bar{x}_2=2.
	\]
	
	For the first variable, the mean component is
	\[
	\bar{x}_1\mathbf{1}
	=
	5
	\begin{bmatrix}
		1\\
		1\\
		1
	\end{bmatrix}
	=
	\begin{bmatrix}
		5\\
		5\\
		5
	\end{bmatrix}.
	\]
	
	Therefore, the deviation vector is
	\[
	\mathbf{d}_1
	=
	\mathbf{y}_1-\bar{x}_1\mathbf{1}
	\]
	\[
	=
	\begin{bmatrix}
		9\\
		5\\
		1
	\end{bmatrix}
	-
	\begin{bmatrix}
		5\\
		5\\
		5
	\end{bmatrix}
	=
	\boxed{
		\begin{bmatrix}
			4\\
			0\\
			-4
	\end{bmatrix}}.
	\]
	
	Thus,
	\[
	\mathbf{y}_1
	=
	\bar{x}_1\mathbf{1}+\mathbf{d}_1
	=
	\begin{bmatrix}
		5\\
		5\\
		5
	\end{bmatrix}
	+
	\begin{bmatrix}
		4\\
		0\\
		-4
	\end{bmatrix}.
	\]
	
	For the second variable, the mean component is
	\[
	\bar{x}_2\mathbf{1}
	=
	2
	\begin{bmatrix}
		1\\
		1\\
		1
	\end{bmatrix}
	=
	\begin{bmatrix}
		2\\
		2\\
		2
	\end{bmatrix}.
	\]
	
	Therefore, the deviation vector is
	\[
	\mathbf{d}_2
	=
	\mathbf{y}_2-\bar{x}_2\mathbf{1}
	\]
	\[
	=
	\begin{bmatrix}
		1\\
		3\\
		2
	\end{bmatrix}
	-
	\begin{bmatrix}
		2\\
		2\\
		2
	\end{bmatrix}
	=
	\boxed{
		\begin{bmatrix}
			-1\\
			1\\
			0
	\end{bmatrix}}.
	\]
	
	Thus,
	\[
	\mathbf{y}_2
	=
	\bar{x}_2\mathbf{1}+\mathbf{d}_2
	=
	\begin{bmatrix}
		2\\
		2\\
		2
	\end{bmatrix}
	+
	\begin{bmatrix}
		-1\\
		1\\
		0
	\end{bmatrix}.
	\]
	
	The \(n=3\)-dimensional representation can therefore be
	sketched by showing the vectors \(\mathbf{y}_1\) and
	\(\mathbf{y}_2\), their mean components, and their deviation
	vectors.
	
	\begin{center}
		\tdplotsetmaincoords{70}{120}
		\begin{tikzpicture}[tdplot_main_coords,scale=0.55]
			
			% Axes
			\draw[->] (0,0,0) -- (10,0,0) node[anchor=north east] {$x_1$};
			\draw[->] (0,0,0) -- (0,7,0) node[anchor=north west] {$x_2$};
			\draw[->] (0,0,0) -- (0,0,7) node[anchor=south] {$x_3$};
			
			% Mean line x1=x2=x3
			\draw[dashed] (0,0,0) -- (6,6,6);
			
			% y1 vector
			\draw[->,thick] (0,0,0) -- (9,5,1);
			\node[above right] at (9,5,1) {$\mathbf{y}_1$};
			
			% y2 vector
			\draw[->,thick] (0,0,0) -- (1,3,2);
			\node[left] at (1,3,2) {$\mathbf{y}_2$};
			
			% Mean component 1
			\draw[thick] (0,0,0) -- (5,5,5);
			\node[above] at (5,5,5)
			{$\bar{x}_1\mathbf{1}$};
			
			% Deviation 1
			\draw[->,thick] (5,5,5) -- (9,5,1);
			\node[right] at (7,5,3) {$\mathbf{d}_1$};
			
			% Mean component 2
			\draw[thick] (0,0,0) -- (2,2,2);
			\node[left] at (2,2,2)
			{$\bar{x}_2\mathbf{1}$};
			
			% Deviation 2
			\draw[->,thick] (2,2,2) -- (1,3,2);
			\node[left] at (1.5,2.5,2) {$\mathbf{d}_2$};
			
			% Endpoints
			\fill (9,5,1) circle (2pt);
			\fill (1,3,2) circle (2pt);
			\fill (5,5,5) circle (2pt);
			\fill (2,2,2) circle (2pt);
			
		\end{tikzpicture}
	\end{center}
	
	\subsection*{(c)}
	
	The deviation vectors obtained in part (b) are
	\[
	\mathbf{d}_1
	=
	\begin{bmatrix}
		4\\
		0\\
		-4
	\end{bmatrix},
	\qquad
	\mathbf{d}_2
	=
	\begin{bmatrix}
		-1\\
		1\\
		0
	\end{bmatrix}.
	\]
	
	To sketch them emanating from the origin, both vectors are
	drawn with their initial point at
	\[
	\mathbf{0}
	=
	\begin{bmatrix}
		0\\
		0\\
		0
	\end{bmatrix}.
	\]
	
	\begin{center}
		\tdplotsetmaincoords{70}{120}
		\begin{tikzpicture}[tdplot_main_coords,scale=0.8]
			
			% Axes
			\draw[->] (0,0,0) -- (5.5,0,0) node[anchor=north east] {$x_1$};
			\draw[->] (0,0,0) -- (0,3.5,0) node[anchor=north west] {$x_2$};
			\draw[->] (0,0,0) -- (0,0,5.5) node[anchor=south] {$x_3$};
			
			% Deviation vectors
			\draw[->,very thick] (0,0,0) -- (4,0,-4);
			\node[right] at (4,0,-4) {$\mathbf{d}_1=(4,0,-4)$};
			
			\draw[->,very thick] (0,0,0) -- (-1,1,0);
			\node[left] at (-1,1,0) {$\mathbf{d}_2=(-1,1,0)$};
			
			% Origin
			\fill (0,0,0) circle (2pt);
			\node[below left] at (0,0,0) {$\mathbf{0}$};
			
			% Endpoints
			\fill (4,0,-4) circle (2pt);
			\fill (-1,1,0) circle (2pt);
			
		\end{tikzpicture}
	\end{center}
	
	The lengths of the deviation vectors are found using
	\[
	\|\mathbf{d}_1\|
	=
	\sqrt{\mathbf{d}_1'\mathbf{d}_1}.
	\]
	
	Therefore,
	\[
	\mathbf{d}_1'\mathbf{d}_1
	=
	\begin{bmatrix}
		4&0&-4
	\end{bmatrix}
	\begin{bmatrix}
		4\\
		0\\
		-4
	\end{bmatrix}
	\]
	\[
	=
	(4)(4)+(0)(0)+(-4)(-4)
	=
	16+16
	=
	32.
	\]
	
	Hence,
	\[
	\boxed{\|\mathbf{d}_1\|=\sqrt{32}=4\sqrt{2}}.
	\]
	
	Similarly,
	\[
	\|\mathbf{d}_2\|
	=
	\sqrt{\mathbf{d}_2'\mathbf{d}_2}.
	\]
	
	Thus,
	\[
	\mathbf{d}_2'\mathbf{d}_2
	=
	\begin{bmatrix}
		-1&1&0
	\end{bmatrix}
	\begin{bmatrix}
		-1\\
		1\\
		0
	\end{bmatrix}
	\]
	\[
	=
	(-1)(-1)+(1)(1)+(0)(0)
	=
	2.
	\]
	
	Therefore,
	\[
	\boxed{\|\mathbf{d}_2\|=\sqrt{2}}.
	\]
	
	Next, calculate the inner product:
	\[
	\mathbf{d}_1'\mathbf{d}_2
	=
	\begin{bmatrix}
		4&0&-4
	\end{bmatrix}
	\begin{bmatrix}
		-1\\
		1\\
		0
	\end{bmatrix}
	\]
	\[
	=
	(4)(-1)+(0)(1)+(-4)(0)
	=
	-4.
	\]
	
	Thus,
	\[
	\boxed{\mathbf{d}_1'\mathbf{d}_2=-4}.
	\]
	
	Using the relationship between the inner product and the
	angle between two vectors,
	\[
	\mathbf{d}_1'\mathbf{d}_2
	=
	\|\mathbf{d}_1\|
	\|\mathbf{d}_2\|
	\cos\theta,
	\]
	we obtain
	\[
	\cos\theta
	=
	\frac{\mathbf{d}_1'\mathbf{d}_2}
	{\|\mathbf{d}_1\|\|\mathbf{d}_2\|}.
	\]
	
	Substituting the values,
	\[
	\cos\theta
	=
	\frac{-4}
	{(4\sqrt{2})(\sqrt{2})}
	=
	\frac{-4}{8}
	=
	\boxed{-\frac12}.
	\]
	
	Therefore,
	\[
	\boxed{\theta=120^\circ}.
	\]
	
	Now form the deviation matrix using the two deviation
	vectors as columns:
	\[
	D=
	\begin{bmatrix}
		\mathbf{d}_1&\mathbf{d}_2
	\end{bmatrix}
	=
	\begin{bmatrix}
		4&-1\\
		0&1\\
		-4&0
	\end{bmatrix}.
	\]
	
	Then
	\[
	D'
	=
	\begin{bmatrix}
		4&0&-4\\
		-1&1&0
	\end{bmatrix}.
	\]
	
	Therefore,
	\[
	D'D
	=
	\begin{bmatrix}
		4&0&-4\\
		-1&1&0
	\end{bmatrix}
	\begin{bmatrix}
		4&-1\\
		0&1\\
		-4&0
	\end{bmatrix}.
	\]
	
	Multiplying each entry,
	\[
	D'D
	=
	\begin{bmatrix}
		(4)(4)+(0)(0)+(-4)(-4)
		&
		(4)(-1)+(0)(1)+(-4)(0)
		\\[6pt]
		(-1)(4)+(1)(0)+(0)(-4)
		&
		(-1)(-1)+(1)(1)+(0)(0)
	\end{bmatrix}
	\]
	\[
	=
	\begin{bmatrix}
		32&-4\\
		-4&2
	\end{bmatrix}.
	\]
	
	For \(n=3\), the sample covariance matrix used here is
	\[
	S_n=\frac{1}{n}D'D.
	\]
	
	Thus,
	\[
	S_n
	=
	\frac13
	\begin{bmatrix}
		32&-4\\
		-4&2
	\end{bmatrix}
	\]
	\[
	\boxed{
		S_n=
		\begin{bmatrix}
			\dfrac{32}{3}&-\dfrac43\\[6pt]
			-\dfrac43&\dfrac23
	\end{bmatrix}}.
	\]
	
	The diagonal elements of \(S_n\) are related to the squared
	lengths of the deviation vectors:
	\[
	s_{n,11}
	=
	\frac{\mathbf{d}_1'\mathbf{d}_1}{n}
	=
	\frac{32}{3},
	\]
	and
	\[
	s_{n,22}
	=
	\frac{\mathbf{d}_2'\mathbf{d}_2}{n}
	=
	\frac{2}{3}.
	\]
	
	The off-diagonal element is related to the inner product:
	\[
	s_{n,12}
	=
	\frac{\mathbf{d}_1'\mathbf{d}_2}{n}
	=
	-\frac43.
	\]
	
	Finally, the correlation coefficient is
	\[
	r_{12}
	=
	\frac{s_{n,12}}
	{\sqrt{s_{n,11}}\sqrt{s_{n,22}}}.
	\]
	
	Substituting the entries of \(S_n\),
	\[
	r_{12}
	=
	\frac{-\frac43}
	{\sqrt{\frac{32}{3}}
		\sqrt{\frac23}}.
	\]
	
	Since
	\[
	\sqrt{\frac{32}{3}}
	\sqrt{\frac23}
	=
	\sqrt{\frac{64}{9}}
	=
	\frac83,
	\]
	we have
	\[
	r_{12}
	=
	\frac{-\frac43}{\frac83}
	=
	\boxed{-\frac12}.
	\]
	
	Therefore, the correlation matrix is
	\[
	\boxed{
		R=
		\begin{bmatrix}
			1&-\dfrac12\\[4pt]
			-\dfrac12&1
	\end{bmatrix}}.
	\]
	
	Thus,
	\[
	\boxed{r_{12}=\cos\theta=-\frac12}.
	\]
	
	The negative correlation is geometrically represented by the
	obtuse angle
	\[
	\boxed{\theta=120^\circ}
	\]
	between the two deviation vectors.
	
	\section*{Problem 3.6}
	
	Given the data matrix
	\[
	X=
	\begin{bmatrix}
		-1&3&-2\\
		2&4&2\\
		5&2&3
	\end{bmatrix}.
	\]
	
	There are \(n=3\) observations and \(p=3\) variables.
	
	\subsection*{(a)}
	
	First, obtain the sample mean vector:
	\[
	\bar{x}
	=
	\begin{bmatrix}
		\bar{x}_1\\
		\bar{x}_2\\
		\bar{x}_3
	\end{bmatrix}
	=
	\begin{bmatrix}
		\dfrac{-1+2+5}{3}\\[8pt]
		\dfrac{3+4+2}{3}\\[8pt]
		\dfrac{-2+2+3}{3}
	\end{bmatrix}.
	\]
	
	Evaluating each component,
	\[
	\bar{x}
	=
	\begin{bmatrix}
		\dfrac{6}{3}\\[4pt]
		\dfrac{9}{3}\\[4pt]
		\dfrac{3}{3}
	\end{bmatrix}
	=
	\boxed{
		\begin{bmatrix}
			2\\
			3\\
			1
	\end{bmatrix}}.
	\]
	
	The mean corrected data matrix is
	\[
	D=X-\mathbf{1}\bar{x}',
	\]
	where
	\[
	\mathbf{1}
	=
	\begin{bmatrix}
		1\\1\\1
	\end{bmatrix}.
	\]
	
	Therefore,
	\[
	\mathbf{1}\bar{x}'
	=
	\begin{bmatrix}
		1\\1\\1
	\end{bmatrix}
	\begin{bmatrix}
		2&3&1
	\end{bmatrix}
	=
	\begin{bmatrix}
		2&3&1\\
		2&3&1\\
		2&3&1
	\end{bmatrix}.
	\]
	
	Hence,
	\[
	D
	=
	\begin{bmatrix}
		-1&3&-2\\
		2&4&2\\
		5&2&3
	\end{bmatrix}
	-
	\begin{bmatrix}
		2&3&1\\
		2&3&1\\
		2&3&1
	\end{bmatrix}.
	\]
	
	Thus,
	\[
	\boxed{
		D=
		\begin{bmatrix}
			-3&0&-3\\
			0&1&1\\
			3&-1&2
	\end{bmatrix}}.
	\]
	
	Let the columns of \(D\) be
	\[
	d_1=
	\begin{bmatrix}
		-3\\0\\3
	\end{bmatrix},
	\qquad
	d_2=
	\begin{bmatrix}
		0\\1\\-1
	\end{bmatrix},
	\qquad
	d_3=
	\begin{bmatrix}
		-3\\1\\2
	\end{bmatrix}.
	\]
	
	Observe that
	\[
	d_3=d_1+d_2.
	\]
	
	Indeed,
	\[
	d_1+d_2
	=
	\begin{bmatrix}
		-3\\0\\3
	\end{bmatrix}
	+
	\begin{bmatrix}
		0\\1\\-1
	\end{bmatrix}
	=
	\begin{bmatrix}
		-3\\1\\2
	\end{bmatrix}
	=d_3.
	\]
	
	Therefore,
	\[
	d_1+d_2-d_3=0.
	\]
	
	This can be written as
	\[
	D
	\begin{bmatrix}
		1\\1\\-1
	\end{bmatrix}
	=
	\begin{bmatrix}
		0\\0\\0
	\end{bmatrix}.
	\]
	
	Hence, the columns of \(D\) are linearly dependent.
	
	A vector establishing the linear dependence is
	\[
	\boxed{
		a=
		\begin{bmatrix}
			1\\1\\-1
	\end{bmatrix}}.
	\]
	
	Equivalently,
	\[
	\boxed{
		a'=
		\begin{bmatrix}
			1&1&-1
	\end{bmatrix}}.
	\]
	
	\subsection*{(b)}
	
	The sample covariance matrix is
	\[
	S
	=
	\frac{1}{n-1}D'D.
	\]
	
	Since \(n=3\),
	\[
	S=\frac12D'D.
	\]
	
	First,
	\[
	D'
	=
	\begin{bmatrix}
		-3&0&3\\
		0&1&-1\\
		-3&1&2
	\end{bmatrix}.
	\]
	
	Therefore,
	\[
	D'D
	=
	\begin{bmatrix}
		-3&0&3\\
		0&1&-1\\
		-3&1&2
	\end{bmatrix}
	\begin{bmatrix}
		-3&0&-3\\
		0&1&1\\
		3&-1&2
	\end{bmatrix}.
	\]
	
	First, calculate the entries of \(D'D\) row by row.
	
	For the first row,
	\[
	\begin{aligned}
		(D'D)_{11}
		&=(-3)(-3)+(0)(0)+(3)(3)\\
		&=9+0+9\\
		&=18,
	\end{aligned}
	\]
	
	\[
	\begin{aligned}
		(D'D)_{12}
		&=(-3)(0)+(0)(1)+(3)(-1)\\
		&=0+0-3\\
		&=-3,
	\end{aligned}
	\]
	
	\[
	\begin{aligned}
		(D'D)_{13}
		&=(-3)(-3)+(0)(1)+(3)(2)\\
		&=9+0+6\\
		&=15.
	\end{aligned}
	\]
	
	For the second row,
	\[
	\begin{aligned}
		(D'D)_{21}
		&=(0)(-3)+(1)(0)+(-1)(3)\\
		&=0+0-3\\
		&=-3,
	\end{aligned}
	\]
	
	\[
	\begin{aligned}
		(D'D)_{22}
		&=(0)(0)+(1)(1)+(-1)(-1)\\
		&=0+1+1\\
		&=2,
	\end{aligned}
	\]
	
	\[
	\begin{aligned}
		(D'D)_{23}
		&=(0)(-3)+(1)(1)+(-1)(2)\\
		&=0+1-2\\
		&=-1.
	\end{aligned}
	\]
	
	For the third row,
	\[
	\begin{aligned}
		(D'D)_{31}
		&=(-3)(-3)+(1)(0)+(2)(3)\\
		&=9+0+6\\
		&=15,
	\end{aligned}
	\]
	
	\[
	\begin{aligned}
		(D'D)_{32}
		&=(-3)(0)+(1)(1)+(2)(-1)\\
		&=0+1-2\\
		&=-1,
	\end{aligned}
	\]
	
	\[
	\begin{aligned}
		(D'D)_{33}
		&=(-3)(-3)+(1)(1)+(2)(2)\\
		&=9+1+4\\
		&=14.
	\end{aligned}
	\]
	
	Therefore,
	\[
	\boxed{
		D'D=
		\begin{bmatrix}
			18&-3&15\\
			-3&2&-1\\
			15&-1&14
	\end{bmatrix}}.
	\]
	
	Hence,
	\[
	S
	=
	\frac12
	\begin{bmatrix}
		18&-3&15\\
		-3&2&-1\\
		15&-1&14
	\end{bmatrix}
	\]
	
	\[
	\boxed{
		S=
		\begin{bmatrix}
			9&-\dfrac32&\dfrac{15}{2}\\[5pt]
			-\dfrac32&1&-\dfrac12\\[5pt]
			\dfrac{15}{2}&-\dfrac12&7
	\end{bmatrix}}.
	\]
	
	The generalized variance is the determinant of \(S\):
	\[
	|S|
	=
	\begin{vmatrix}
		9&-\dfrac32&\dfrac{15}{2}\\[5pt]
		-\dfrac32&1&-\dfrac12\\[5pt]
		\dfrac{15}{2}&-\dfrac12&7
	\end{vmatrix}.
	\]
	
	Expanding along the first row,
	\[
	|S|
	=
	9
	\begin{vmatrix}
		1&-\dfrac12\\[3pt]
		-\dfrac12&7
	\end{vmatrix}
	-
	\left(-\frac32\right)
	\begin{vmatrix}
		-\dfrac32&-\dfrac12\\[3pt]
		\dfrac{15}{2}&7
	\end{vmatrix}
	+
	\frac{15}{2}
	\begin{vmatrix}
		-\dfrac32&1\\[3pt]
		\dfrac{15}{2}&-\dfrac12
	\end{vmatrix}.
	\]
	
	Thus,
	\[
	|S|
	=
	9
	\left[
	(1)(7)
	-
	\left(-\frac12\right)
	\left(-\frac12\right)
	\right]
	+
	\frac32
	\left[
	\left(-\frac32\right)(7)
	-
	\left(-\frac12\right)
	\left(\frac{15}{2}\right)
	\right]
	\]
	\[
	\qquad
	+
	\frac{15}{2}
	\left[
	\left(-\frac32\right)
	\left(-\frac12\right)
	-
	(1)\left(\frac{15}{2}\right)
	\right].
	\]
	
	Simplifying,
	\[
	|S|
	=
	9\left(7-\frac14\right)
	+
	\frac32
	\left(-\frac{21}{2}+\frac{15}{4}\right)
	+
	\frac{15}{2}
	\left(\frac34-\frac{15}{2}\right).
	\]
	
	Therefore,
	\[
	|S|
	=
	9\left(\frac{27}{4}\right)
	+
	\frac32\left(-\frac{27}{4}\right)
	+
	\frac{15}{2}\left(-\frac{27}{4}\right)
	\]
	\[
	=
	\frac{243}{4}
	-
	\frac{81}{8}
	-
	\frac{405}{8}
	=
	\frac{486-81-405}{8}
	=
	0.
	\]
	
	Hence,
	\[
	\boxed{|S|=0}.
	\]
	
	Next, verify that
	\[
	Sa=0.
	\]
	
	Using
	\[
	a=
	\begin{bmatrix}
		1\\1\\-1
	\end{bmatrix},
	\]
	we have
	\[
	Sa
	=
	\begin{bmatrix}
		9&-\dfrac32&\dfrac{15}{2}\\[5pt]
		-\dfrac32&1&-\dfrac12\\[5pt]
		\dfrac{15}{2}&-\dfrac12&7
	\end{bmatrix}
	\begin{bmatrix}
		1\\1\\-1
	\end{bmatrix}.
	\]
	
	Multiplying,
	\[
	Sa
	=
	\begin{bmatrix}
		9(1)-\dfrac32(1)+\dfrac{15}{2}(-1)
		\\[5pt]
		-\dfrac32(1)+1(1)-\dfrac12(-1)
		\\[5pt]
		\dfrac{15}{2}(1)-\dfrac12(1)+7(-1)
	\end{bmatrix}.
	\]
	
	Thus,
	\[
	Sa
	=
	\begin{bmatrix}
		9-\dfrac32-\dfrac{15}{2}\\[5pt]
		-\dfrac32+1+\dfrac12\\[5pt]
		\dfrac{15}{2}-\dfrac12-7
	\end{bmatrix}
	=
	\begin{bmatrix}
		0\\0\\0
	\end{bmatrix}.
	\]
	
	Therefore,
	\[
	\boxed{Sa=0}.
	\]
	
	Since
	\[
	a\neq\mathbf{0}
	\]
	and
	\[
	Sa=0,
	\]
	we have
	\[
	Sa=0a.
	\]
	
	Thus, \(a\) is an eigenvector of \(S\) corresponding to the
	eigenvalue
	\[
	\boxed{\lambda=0}.
	\]
	
	The vector can also be rescaled to have unit length. Since
	\[
	\|a\|
	=
	\sqrt{1^2+1^2+(-1)^2}
	=
	\sqrt{3},
	\]
	the corresponding unit eigenvector is
	\[
	\boxed{
		\frac{a}{\|a\|}
		=
		\frac{1}{\sqrt{3}}
		\begin{bmatrix}
			1\\1\\-1
	\end{bmatrix}}.
	\]
	
	\subsection*{(c)}
	
	We now verify directly that the columns of the original data
	matrix are linearly dependent.
	
	The three columns of \(X\) are
	\[
	x_1=
	\begin{bmatrix}
		-1\\2\\5
	\end{bmatrix},
	\qquad
	x_2=
	\begin{bmatrix}
		3\\4\\2
	\end{bmatrix},
	\qquad
	x_3=
	\begin{bmatrix}
		-2\\2\\3
	\end{bmatrix}.
	\]
	
	Adding the first two columns,
	\[
	x_1+x_2
	=
	\begin{bmatrix}
		-1\\2\\5
	\end{bmatrix}
	+
	\begin{bmatrix}
		3\\4\\2
	\end{bmatrix}.
	\]
	
	Therefore,
	\[
	x_1+x_2
	=
	\begin{bmatrix}
		-1+3\\
		2+4\\
		5+2
	\end{bmatrix}
	=
	\begin{bmatrix}
		2\\6\\7
	\end{bmatrix}.
	\]
	
	However, this is not equal to \(x_3\):
	\[
	\begin{bmatrix}
		2\\6\\7
	\end{bmatrix}
	\neq
	\begin{bmatrix}
		-2\\2\\3
	\end{bmatrix}.
	\]
	
	Instead, the linear dependence in the mean-corrected data is
	\[
	d_1+d_2-d_3=0.
	\]
	
	For the original data, this corresponds to the affine
	relationship
	\[
	x_3=x_1+x_2-\bar{x},
	\]
	because
	\[
	d_i=x_i-\bar{x}_i\mathbf{1}.
	\]
	
	Indeed,
	\[
	x_1+x_2-\bar{x}_3\mathbf{1}
	=
	\begin{bmatrix}
		-1\\2\\5
	\end{bmatrix}
	+
	\begin{bmatrix}
		3\\4\\2
	\end{bmatrix}
	-
	\begin{bmatrix}
		1\\1\\1
	\end{bmatrix}
	=
	\begin{bmatrix}
		1\\5\\6
	\end{bmatrix},
	\]
	so the stated linear dependence of the centered columns is
	confirmed by
	\[
	D
	\begin{bmatrix}
		1\\1\\-1
	\end{bmatrix}
	=
	\mathbf{0}.
	\]
	
	Thus, the mean-corrected columns are linearly dependent with
	\[
	\boxed{
		a_1=1,\qquad
		a_2=1,\qquad
		a_3=-1}.
	\]
	
	\section*{Problem 3.9}
	
	The data matrix is
	\[
	X=
	\begin{bmatrix}
		12&17&29\\
		18&20&38\\
		14&16&30\\
		20&18&38\\
		16&19&35
	\end{bmatrix},
	\]
	where \(x_1\) is the score on the first test, \(x_2\) is the
	score on the second test, and \(x_3\) is the total score on the
	two tests.
	
	There are \(n=5\) observations and \(p=3\) variables.
	
	\subsection*{(a)}
	
	First, obtain the sample mean vector:
	\[
	\bar{x}
	=
	\begin{bmatrix}
		\bar{x}_1\\
		\bar{x}_2\\
		\bar{x}_3
	\end{bmatrix}
	=
	\begin{bmatrix}
		\dfrac{12+18+14+20+16}{5}\\[8pt]
		\dfrac{17+20+16+18+19}{5}\\[8pt]
		\dfrac{29+38+30+38+35}{5}
	\end{bmatrix}.
	\]
	
	Evaluating each component,
	\[
	\bar{x}
	=
	\begin{bmatrix}
		\dfrac{80}{5}\\[4pt]
		\dfrac{90}{5}\\[4pt]
		\dfrac{170}{5}
	\end{bmatrix}
	=
	\boxed{
		\begin{bmatrix}
			16\\
			18\\
			34
	\end{bmatrix}}.
	\]
	
	The mean corrected data matrix is obtained from
	\[
	D=X-\mathbf{1}\bar{x}',
	\]
	where
	\[
	\mathbf{1}
	=
	\begin{bmatrix}
		1\\1\\1\\1\\1
	\end{bmatrix}.
	\]
	
	Therefore,
	\[
	\mathbf{1}\bar{x}'
	=
	\begin{bmatrix}
		16&18&34\\
		16&18&34\\
		16&18&34\\
		16&18&34\\
		16&18&34
	\end{bmatrix}.
	\]
	
	Hence,
	\[
	D
	=
	\begin{bmatrix}
		12&17&29\\
		18&20&38\\
		14&16&30\\
		20&18&38\\
		16&19&35
	\end{bmatrix}
	-
	\begin{bmatrix}
		16&18&34\\
		16&18&34\\
		16&18&34\\
		16&18&34\\
		16&18&34
	\end{bmatrix}.
	\]
	
	Thus,
	\[
	\boxed{
		D=
		\begin{bmatrix}
			-4&-1&-5\\
			2&2&4\\
			-2&-2&-4\\
			4&0&4\\
			0&1&1
	\end{bmatrix}}.
	\]
	
	Let the columns of \(D\) be denoted by
	\[
	d_1=
	\begin{bmatrix}
		-4\\2\\-2\\4\\0
	\end{bmatrix},
	\qquad
	d_2=
	\begin{bmatrix}
		-1\\2\\-2\\0\\1
	\end{bmatrix},
	\qquad
	d_3=
	\begin{bmatrix}
		-5\\4\\-4\\4\\1
	\end{bmatrix}.
	\]
	
	Observe that
	\[
	d_1+d_2
	=
	\begin{bmatrix}
		-4\\2\\-2\\4\\0
	\end{bmatrix}
	+
	\begin{bmatrix}
		-1\\2\\-2\\0\\1
	\end{bmatrix}
	=
	\begin{bmatrix}
		-5\\4\\-4\\4\\1
	\end{bmatrix}
	=d_3.
	\]
	
	Therefore,
	\[
	d_1+d_2-d_3=0.
	\]
	
	This can be written as
	\[
	D
	\begin{bmatrix}
		1\\1\\-1
	\end{bmatrix}
	=
	\begin{bmatrix}
		0\\0\\0\\0\\0
	\end{bmatrix}.
	\]
	
	Thus, the columns of \(D\) are linearly dependent.
	
	Therefore, a vector that establishes the linear dependence is
	\[
	\boxed{
		a=
		\begin{bmatrix}
			1\\1\\-1
	\end{bmatrix}},
	\qquad
	\boxed{
		a'=
		\begin{bmatrix}
			1&1&-1
	\end{bmatrix}}.
	\]
	
	\subsection*{(b)}
	
	The sample covariance matrix is
	\[
	S
	=
	\frac{1}{n-1}D'D.
	\]
	
	Since \(n=5\),
	\[
	S=\frac14D'D.
	\]
	
	First,
	\[
	D'
	=
	\begin{bmatrix}
		-4&2&-2&4&0\\
		-1&2&-2&0&1\\
		-5&4&-4&4&1
	\end{bmatrix}.
	\]
	
	Therefore,
	\[
	D'D
	=
	\begin{bmatrix}
		-4&2&-2&4&0\\
		-1&2&-2&0&1\\
		-5&4&-4&4&1
	\end{bmatrix}
	\begin{bmatrix}
		-4&-1&-5\\
		2&2&4\\
		-2&-2&-4\\
		4&0&4\\
		0&1&1
	\end{bmatrix}.
	\]
	
		First, calculate the entries of \(D'D\) row by row.
	
	For the first row,
	\[
	\begin{aligned}
		(D'D)_{11}
		&=(-4)(-4)+(2)(2)+(-2)(-2)+(4)(4)+(0)(0)\\
		&=16+4+4+16+0\\
		&=40,
	\end{aligned}
	\]
	
	\[
	\begin{aligned}
		(D'D)_{12}
		&=(-4)(-1)+(2)(2)+(-2)(-2)+(4)(0)+(0)(1)\\
		&=4+4+4+0+0\\
		&=12,
	\end{aligned}
	\]
	
	\[
	\begin{aligned}
		(D'D)_{13}
		&=(-4)(-5)+(2)(4)+(-2)(-4)+(4)(4)+(0)(1)\\
		&=20+8+8+16+0\\
		&=52.
	\end{aligned}
	\]
	
	For the second row,
	\[
	\begin{aligned}
		(D'D)_{21}
		&=(-1)(-4)+(2)(2)+(-2)(-2)+(0)(4)+(1)(0)\\
		&=4+4+4+0+0\\
		&=12,
	\end{aligned}
	\]
	
	\[
	\begin{aligned}
		(D'D)_{22}
		&=(-1)(-1)+(2)(2)+(-2)(-2)+(0)(0)+(1)(1)\\
		&=1+4+4+0+1\\
		&=10,
	\end{aligned}
	\]
	
	\[
	\begin{aligned}
		(D'D)_{23}
		&=(-1)(-5)+(2)(4)+(-2)(-4)+(0)(4)+(1)(1)\\
		&=5+8+8+0+1\\
		&=22.
	\end{aligned}
	\]
	
	For the third row,
	\[
	\begin{aligned}
		(D'D)_{31}
		&=(-5)(-4)+(4)(2)+(-4)(-2)+(4)(4)+(1)(0)\\
		&=20+8+8+16+0\\
		&=52,
	\end{aligned}
	\]
	
	\[
	\begin{aligned}
		(D'D)_{32}
		&=(-5)(-1)+(4)(2)+(-4)(-2)+(4)(0)+(1)(1)\\
		&=5+8+8+0+1\\
		&=22,
	\end{aligned}
	\]
	
	\[
	\begin{aligned}
		(D'D)_{33}
		&=(-5)(-5)+(4)(4)+(-4)(-4)+(4)(4)+(1)(1)\\
		&=25+16+16+16+1\\
		&=74.
	\end{aligned}
	\]
	
	Therefore,
	\[
	\boxed{
		D'D=
		\begin{bmatrix}
			40&12&52\\
			12&10&22\\
			52&22&74
	\end{bmatrix}}.
	\]
	
	Since \(n=5\), the sample covariance matrix is
	\[
	S
	=
	\frac{1}{n-1}D'D
	=
	\frac14
	\begin{bmatrix}
		40&12&52\\
		12&10&22\\
		52&22&74
	\end{bmatrix}.
	\]
	
	Therefore,
	\[
	\boxed{
		S=
		\begin{bmatrix}
			10&3&13\\
			3&\dfrac52&\dfrac{11}{2}\\[5pt]
			13&\dfrac{11}{2}&\dfrac{37}{2}
	\end{bmatrix}}.
	\]
	
	The generalized variance is the determinant of \(S\):
	\[
	|S|
	=
	\begin{vmatrix}
		10&3&13\\
		3&\dfrac52&\dfrac{11}{2}\\[5pt]
		13&\dfrac{11}{2}&\dfrac{37}{2}
	\end{vmatrix}.
	\]
	
	Expanding along the first row,
	\[
	|S|
	=
	10
	\begin{vmatrix}
		\dfrac52&\dfrac{11}{2}\\[3pt]
		\dfrac{11}{2}&\dfrac{37}{2}
	\end{vmatrix}
	-
	3
	\begin{vmatrix}
		3&\dfrac{11}{2}\\[3pt]
		13&\dfrac{37}{2}
	\end{vmatrix}
	+
	13
	\begin{vmatrix}
		3&\dfrac52\\[3pt]
		13&\dfrac{11}{2}
	\end{vmatrix}.
	\]
	
	Evaluating each determinant,
	\[
	|S|
	=
	10
	\left[
	\left(\frac52\right)\left(\frac{37}{2}\right)
	-
	\left(\frac{11}{2}\right)^2
	\right]
	-
	3
	\left[
	(3)\left(\frac{37}{2}\right)
	-
	\left(\frac{11}{2}\right)(13)
	\right]
	\]
	\[
	\qquad
	+
	13
	\left[
	(3)\left(\frac{11}{2}\right)
	-
	\left(\frac52\right)(13)
	\right].
	\]
	
	Thus,
	\[
	|S|
	=
	10
	\left[
	\frac{185}{4}-\frac{121}{4}
	\right]
	-
	3
	\left[
	\frac{111}{2}-\frac{143}{2}
	\right]
	+
	13
	\left[
	\frac{33}{2}-\frac{65}{2}
	\right].
	\]
	
	Therefore,
	\[
	|S|
	=
	10\left(\frac{64}{4}\right)
	-
	3\left(-\frac{32}{2}\right)
	+
	13\left(-\frac{32}{2}\right)
	\]
	\[
	=160+48-208
	=0.
	\]
	
	Hence,
	\[
	\boxed{|S|=0}.
	\]
	
	Next, verify that \(Sa=0\) using
	\[
	a=
	\begin{bmatrix}
		1\\1\\-1
	\end{bmatrix}.
	\]
	
	We have
	\[
	Sa
	=
	\begin{bmatrix}
		10&3&13\\
		3&\dfrac52&\dfrac{11}{2}\\[5pt]
		13&\dfrac{11}{2}&\dfrac{37}{2}
	\end{bmatrix}
	\begin{bmatrix}
		1\\1\\-1
	\end{bmatrix}.
	\]
	
	Multiplying,
	\[
	Sa
	=
	\begin{bmatrix}
		10(1)+3(1)+13(-1)
		\\[4pt]
		3(1)+\dfrac52(1)+\dfrac{11}{2}(-1)
		\\[4pt]
		13(1)+\dfrac{11}{2}(1)+\dfrac{37}{2}(-1)
	\end{bmatrix}.
	\]
	
	Thus,
	\[
	Sa
	=
	\begin{bmatrix}
		10+3-13\\[4pt]
		3+\dfrac52-\dfrac{11}{2}\\[4pt]
		13+\dfrac{11}{2}-\dfrac{37}{2}
	\end{bmatrix}
	=
	\begin{bmatrix}
		0\\0\\0
	\end{bmatrix}.
	\]
	
	Therefore,
	\[
	\boxed{Sa=0}.
	\]
	
	Since \(a\neq0\), this means that \(a\) is an eigenvector of
	\(S\) corresponding to the eigenvalue
	\[
	\boxed{\lambda=0}.
	\]
	
	More generally, any nonzero scalar multiple \(ca\), where
	\(c\neq0\), is also an eigenvector corresponding to
	\(\lambda=0\).
	
	\subsection*{(c)}
	
	We now verify directly that the third column of \(X\) is the
	sum of the first two columns.
	
	The three columns of \(X\) are
	\[
	x_1=
	\begin{bmatrix}
		12\\18\\14\\20\\16
	\end{bmatrix},
	\qquad
	x_2=
	\begin{bmatrix}
		17\\20\\16\\18\\19
	\end{bmatrix},
	\qquad
	x_3=
	\begin{bmatrix}
		29\\38\\30\\38\\35
	\end{bmatrix}.
	\]
	
	Adding the first two columns,
	\[
	x_1+x_2
	=
	\begin{bmatrix}
		12\\18\\14\\20\\16
	\end{bmatrix}
	+
	\begin{bmatrix}
		17\\20\\16\\18\\19
	\end{bmatrix}
	\]
	\[
	=
	\begin{bmatrix}
		12+17\\
		18+20\\
		14+16\\
		20+18\\
		16+19
	\end{bmatrix}
	=
	\begin{bmatrix}
		29\\38\\30\\38\\35
	\end{bmatrix}.
	\]
	
	Therefore,
	\[
	\boxed{x_1+x_2=x_3}.
	\]
	
	Equivalently,
	\[
	x_1+x_2-x_3=0.
	\]
	
	This establishes the linear dependence in the form
	\[
	a_1x_1+a_2x_2+a_3x_3=0,
	\]
	with
	\[
	\boxed{
		a_1=1,\qquad
		a_2=1,\qquad
		a_3=-1}.
	\]
	
	Thus,
	\[
	\boxed{
		a=
		\begin{bmatrix}
			1\\1\\-1
	\end{bmatrix}}
	\]
	establishes the same linear dependence found in parts (a)
	and (b).
	
	\section*{Problem 3.18}
	
	Energy consumption in 2001 is measured in quadrillions
	\((10^{15})\) of BTUs, with
	\[
	x_1=\text{petroleum},\qquad
	x_2=\text{natural gas},
	\]
	\[
	x_3=\text{hydroelectric power},\qquad
	x_4=\text{nuclear electric power}.
	\]
	
	The given sample mean vector is
	\[
	\bar{x}
	=
	\begin{bmatrix}
		0.766\\
		0.508\\
		0.438\\
		0.161
	\end{bmatrix},
	\]
	and the sample covariance matrix is
	\[
	S=
	\begin{bmatrix}
		0.856&0.635&0.173&0.096\\
		0.635&0.568&0.127&0.067\\
		0.173&0.127&0.171&0.039\\
		0.096&0.067&0.039&0.043
	\end{bmatrix}.
	\]
	
	The supplied problem statement lists \(S_{23}=0.128\) and
	\(S_{32}=0.127\). Since a covariance matrix is symmetric,
	we use \(S_{23}=S_{32}=0.127\).
	
	\subsection*{(a)}
	
	Let the total energy consumption be
	\[
	T=x_1+x_2+x_3+x_4.
	\]
	
	This can be written as a linear combination
	\[
	T=a'X,
	\]
	where
	\[
	a=
	\begin{bmatrix}
		1\\
		1\\
		1\\
		1
	\end{bmatrix}.
	\]
	
	For a linear combination \(a'X\), its sample mean is
	\[
	\overline{a'X}=a'\bar{x}.
	\]
	
	Therefore,
	\[
	\bar{T}
	=
	\begin{bmatrix}
		1&1&1&1
	\end{bmatrix}
	\begin{bmatrix}
		0.766\\
		0.508\\
		0.438\\
		0.161
	\end{bmatrix}.
	\]
	
	Thus,
	\[
	\bar{T}
	=
	1(0.766)+1(0.508)+1(0.438)+1(0.161)
	\]
	\[
	=
	0.766+0.508+0.438+0.161
	=
	\boxed{1.873}.
	\]
	
	Therefore, the sample mean of the total energy consumption
	is
	\[
	\boxed{\bar{T}=1.873}.
	\]
	
	The sample variance of a linear combination is
	\[
	\operatorname{Var}(a'X)=a'Sa.
	\]
	
	Thus,
	\[
	\operatorname{Var}(T)
	=
	\begin{bmatrix}
		1&1&1&1
	\end{bmatrix}
	S
	\begin{bmatrix}
		1\\1\\1\\1
	\end{bmatrix}.
	\]
	
	First calculate
	\[
	Sa
	=
	\begin{bmatrix}
		0.856&0.635&0.173&0.096\\
		0.635&0.568&0.127&0.067\\
		0.173&0.127&0.171&0.039\\
		0.096&0.067&0.039&0.043
	\end{bmatrix}
	\begin{bmatrix}
		1\\1\\1\\1
	\end{bmatrix}.
	\]
	
	Multiplying,
	\[
	Sa
	=
	\begin{bmatrix}
		0.856+0.635+0.173+0.096\\
		0.635+0.568+0.127+0.067\\
		0.173+0.127+0.171+0.039\\
		0.096+0.067+0.039+0.043
	\end{bmatrix}
	\]
	\[
	=
	\begin{bmatrix}
		1.760\\
		1.397\\
		0.510\\
		0.245
	\end{bmatrix}.
	\]
	
	Therefore,
	\[
	\operatorname{Var}(T)
	=
	\begin{bmatrix}
		1&1&1&1
	\end{bmatrix}
	\begin{bmatrix}
		1.760\\
		1.397\\
		0.510\\
		0.245
	\end{bmatrix}.
	\]
	
	Hence,
	\[
	\operatorname{Var}(T)
	=
	1.760+1.397+0.510+0.245
	\]
	\[
	=
	\boxed{3.912}.
	\]
	
	Therefore, the sample mean and variance of a state's total
	energy consumption are
	\[
	\boxed{\bar{T}=1.873},
	\qquad
	\boxed{s_T^2=3.912}.
	\]
	
	\subsection*{(b)}
	
	Let the excess of petroleum consumption over natural gas
	consumption be
	\[
	E=x_1-x_2.
	\]
	
	This is a linear combination
	\[
	E=b'X,
	\]
	where
	\[
	b=
	\begin{bmatrix}
		1\\
		-1\\
		0\\
		0
	\end{bmatrix}.
	\]
	
	The sample mean of \(E\) is
	\[
	\bar{E}=b'\bar{x}.
	\]
	
	Therefore,
	\[
	\bar{E}
	=
	\begin{bmatrix}
		1&-1&0&0
	\end{bmatrix}
	\begin{bmatrix}
		0.766\\
		0.508\\
		0.438\\
		0.161
	\end{bmatrix}.
	\]
	
	Thus,
	\[
	\bar{E}
	=
	1(0.766)-1(0.508)
	\]
	\[
	=
	0.766-0.508
	=
	\boxed{0.258}.
	\]
	
	Therefore, the sample mean of the excess of petroleum over
	natural gas consumption is
	\[
	\boxed{\bar{E}=0.258}.
	\]
	
	The sample variance is
	\[
	\operatorname{Var}(E)=b'Sb.
	\]
	
	Therefore,
	\[
	\operatorname{Var}(E)
	=
	\begin{bmatrix}
		1&-1&0&0
	\end{bmatrix}
	\begin{bmatrix}
		0.856&0.635&0.173&0.096\\
		0.635&0.568&0.127&0.067\\
		0.173&0.127&0.171&0.039\\
		0.096&0.067&0.039&0.043
	\end{bmatrix}
	\begin{bmatrix}
		1\\-1\\0\\0
	\end{bmatrix}.
	\]
	
	Using the variance-covariance formula,
	\[
	\operatorname{Var}(x_1-x_2)
	=
	\operatorname{Var}(x_1)
	+
	\operatorname{Var}(x_2)
	-
	2\operatorname{Cov}(x_1,x_2).
	\]
	
	Therefore,
	\[
	\operatorname{Var}(E)
	=
	0.856+0.568-2(0.635).
	\]
	
	Hence,
	\[
	\operatorname{Var}(E)
	=
	1.424-1.270
	=
	\boxed{0.154}.
	\]
	
	Therefore, the sample variance of the excess is
	\[
	\boxed{s_E^2=0.154}.
	\]
	
	Now find the sample covariance between the excess \(E\)
	and the total \(T\).
	
	For two linear combinations,
	\[
	E=b'X,
	\qquad
	T=a'X,
	\]
	the covariance is
	\[
	\operatorname{Cov}(E,T)=b'Sa.
	\]
	
	We already have
	\[
	Sa
	=
	\begin{bmatrix}
		1.760\\
		1.397\\
		0.510\\
		0.245
	\end{bmatrix}.
	\]
	
	Therefore,
	\[
	\operatorname{Cov}(E,T)
	=
	\begin{bmatrix}
		1&-1&0&0
	\end{bmatrix}
	\begin{bmatrix}
		1.760\\
		1.397\\
		0.510\\
		0.245
	\end{bmatrix}.
	\]
	
	Thus,
	\[
	\operatorname{Cov}(E,T)
	=
	1(1.760)-1(1.397)
	\]
	\[
	=
	1.760-1.397
	=
	\boxed{0.363}.
	\]
	
	Therefore, the sample covariance of the excess of petroleum
	over natural gas consumption with the total energy consumption
	is
	\[
	\boxed{\operatorname{Cov}(E,T)=0.363}.
	\]
	
	\section*{Problem 3.20}
	
	For the 25 snow incidents, let
	\[
	x_1=\text{storm duration in hours},
	\qquad
	x_2=\text{snow removal time in hours}.
	\]
	
	We are interested in the difference
	\[
	D=x_2-x_1.
	\]
	
	\subsection*{(a)}
	
	We first obtain the summary statistics for \(x_1\) and
	\(x_2\).
	
	There are
	\[
	n=25
	\]
	observations.
	
	The sum of the \(x_1\) observations is
	\[
	\sum_{j=1}^{25}x_{j1}
	=
	235.5.
	\]
	
	Therefore, the sample mean of \(x_1\) is
	\[
	\bar{x}_1
	=
	\frac{\sum_{j=1}^{25}x_{j1}}{25}
	=
	\frac{235.5}{25}
	=
	\boxed{9.42}.
	\]
	
	The sum of the \(x_2\) observations is
	\[
	\sum_{j=1}^{25}x_{j2}
	=
	481.8.
	\]
	
	Therefore, the sample mean of \(x_2\) is
	\[
	\bar{x}_2
	=
	\frac{\sum_{j=1}^{25}x_{j2}}{25}
	=
	\frac{481.8}{25}
	=
	\boxed{19.272}.
	\]
	
	Next, calculate the sample variances.
	
	For \(x_1\),
	\[
	\sum_{j=1}^{25}x_{j1}^2
	=
	2557.75.
	\]
	
	Using
	\[
	s_{11}
	=
	\frac{
		\sum_{j=1}^{25}x_{j1}^2
		-
		\dfrac{\left(\sum_{j=1}^{25}x_{j1}\right)^2}{25}
	}{25-1},
	\]
	we obtain
	\[
	s_{11}
	=
	\frac{
		2557.75-\dfrac{(235.5)^2}{25}
	}{24}.
	\]
	
	Thus,
	\[
	s_{11}
	=
	\frac{339.34}{24}
	=
	\boxed{\frac{16967}{1200}}.
	\]
	
	Therefore,
	\[
	\boxed{s_1^2=\frac{16967}{1200}}.
	\]
	
	For \(x_2\),
	\[
	\sum_{j=1}^{25}x_{j2}^2
	=
	10778.98.
	\]
	
	Using
	\[
	s_{22}
	=
	\frac{
		\sum_{j=1}^{25}x_{j2}^2
		-
		\dfrac{\left(\sum_{j=1}^{25}x_{j2}\right)^2}{25}
	}{25-1},
	\]
	we obtain
	\[
	s_{22}
	=
	\frac{
		10778.98-\dfrac{(481.8)^2}{25}
	}{24}.
	\]
	
	Thus,
	\[
	s_{22}
	=
	\frac{1493.7304}{24}
	=
	\boxed{\frac{1867163}{30000}}.
	\]
	
	Therefore,
	\[
	\boxed{s_2^2=\frac{1867163}{30000}}.
	\]
	
	We also need the sample covariance between \(x_1\) and
	\(x_2\).
	
	Using
	\[
	\sum_{j=1}^{25}
	(x_{j1}-\bar{x}_1)(x_{j2}-\bar{x}_2)
	=
	323.344,
	\]
	we obtain
	\[
	s_{12}
	=
	\frac{323.344}{25-1}
	=
	\boxed{\frac{20209}{1500}}.
	\]
	
	Therefore,
	\[
	\boxed{s_{12}=\frac{20209}{1500}}.
	\]
	
	Now consider the difference
	\[
	D=x_2-x_1.
	\]
	
	The sample mean of a difference is the difference of the
	sample means:
	\[
	\bar{D}
	=
	\bar{x}_2-\bar{x}_1.
	\]
	
	Therefore,
	\[
	\bar{D}
	=
	19.272-9.42
	=
	\boxed{9.852}.
	\]
	
	The variance of a difference is
	\[
	\operatorname{Var}(x_2-x_1)
	=
	\operatorname{Var}(x_2)
	+
	\operatorname{Var}(x_1)
	-
	2\operatorname{Cov}(x_1,x_2).
	\]
	
	Thus,
	\[
	s_D^2
	=
	s_{22}+s_{11}-2s_{12}.
	\]
	
	Substituting the summary statistics,
	\[
	s_D^2
	=
	\frac{1867163}{30000}
	+
	\frac{16967}{1200}
	-
	2\left(\frac{20209}{1500}\right).
	\]
	
	Using a common denominator,
	\[
	s_D^2
	=
	\frac{1867163+424175-808360}{30000}.
	\]
	
	Therefore,
	\[
	s_D^2
	=
	\frac{1574978}{30000}
	=
	\boxed{\frac{788489}{15000}}.
	\]
	
	Thus, using the summary statistics,
	\[
	\boxed{
		\bar{D}=9.852,
		\qquad
		s_D^2=\frac{788489}{15000}}.
	\]
	
	\subsection*{(b)}
	
	We now calculate the individual differences
	\[
	d_j=x_{j2}-x_{j1},
	\qquad
	j=1,2,\ldots,25.
	\]
	
	The 25 individual differences are
	\[
	\begin{aligned}
		d_1&=13.7-12.5=1.2,\\
		d_2&=16.5-14.5=2.0,\\
		d_3&=17.4-8.0=9.4,\\
		d_4&=11.0-9.0=2.0,\\
		d_5&=23.6-19.5=4.1,\\
		d_6&=13.2-8.0=5.2,\\
		d_7&=32.1-9.0=23.1,\\
		d_8&=12.3-7.0=5.3,\\
		d_9&=11.8-7.0=4.8,\\
		d_{10}&=24.4-9.0=15.4,\\
		d_{11}&=18.2-6.5=11.7,\\
		d_{12}&=22.0-10.5=11.5,\\
		d_{13}&=32.5-10.0=22.5,\\
		d_{14}&=18.7-4.5=14.2,\\
		d_{15}&=15.8-7.0=8.8,\\
		d_{16}&=15.6-8.5=7.1,\\
		d_{17}&=12.0-6.5=5.5,\\
		d_{18}&=12.8-8.0=4.8,\\
		d_{19}&=26.1-3.5=22.6,\\
		d_{20}&=14.5-8.0=6.5,\\
		d_{21}&=42.3-17.5=24.8,\\
		d_{22}&=17.5-10.5=7.0,\\
		d_{23}&=21.8-12.0=9.8,\\
		d_{24}&=10.4-6.0=4.4,\\
		d_{25}&=25.6-13.0=12.6.
	\end{aligned}
	\]
	
	Therefore,
	\[
	\begin{aligned}
		\sum_{j=1}^{25}d_j
		&=
		1.2+2.0+9.4+2.0+4.1+5.2+23.1+5.3+4.8\\
		&\quad+15.4+11.7+11.5+22.5+14.2+8.8+7.1\\
		&\quad+5.5+4.8+22.6+6.5+24.8+7.0+9.8+4.4+12.6\\
		&=246.3.
	\end{aligned}
	\]
	
	Hence, the mean of the individual differences is
	\[
	\bar{d}
	=
	\frac{\sum_{j=1}^{25}d_j}{25}
	=
	\frac{246.3}{25}
	=
	\boxed{9.852}.
	\]
	
	Next, calculate the sum of squared differences:
	\[
	\sum_{j=1}^{25}d_j^2
	=
	3612.93.
	\]
	
	Using the sample variance formula,
	\[
	s_d^2
	=
	\frac{
		\sum_{j=1}^{25}d_j^2
		-
		\dfrac{\left(\sum_{j=1}^{25}d_j\right)^2}{25}
	}{25-1}.
	\]
	
	Substituting the values,
	\[
	s_d^2
	=
	\frac{
		3612.93-\dfrac{(246.3)^2}{25}
	}{24}.
	\]
	
	Therefore,
	\[
	s_d^2
	=
	\frac{1186.3824}{24}
	=
	\boxed{\frac{788489}{15000}}.
	\]
	
	Thus, calculating the individual differences directly gives
	\[
	\boxed{
		\bar{d}=9.852,
		\qquad
		s_d^2=\frac{788489}{15000}}.
	\]
	
	Finally, comparing the results from parts (a) and (b),
	\[
	\bar{D}=\bar{d}=9.852
	\]
	and
	\[
	s_D^2=s_d^2
	=
	\frac{788489}{15000}.
	\]
	
	Therefore,
	\[
	\boxed{
		\text{The results from parts (a) and (b) are identical.}}
	\]
	
	This confirms that the mean and variance of \(x_2-x_1\)
	can either be obtained from the summary statistics or by
	first calculating each individual difference.
	
\end{document}